\section{Lean formalization architecture}
\label{sec:lean}

\subsection{Imports and namespace}

The reusable core should live beside \code{IPRADerivations}, not inside one
proof packet.  A minimal first import closure is

\begin{lstlisting}
import IPRADerivations.Core.Search
import IPRADerivations.Core.Pairing
import IPRADerivations.Core.Definitions

namespace LeanProofs.IPRA.SCOPE
open LeanProofs.ConciseIPRA LeanProofs.IPRA Syn
open scoped Fm
\end{lstlisting}

The generic library should not import Foundation, Busy Beaver, Diophantine
Consistency, or theorem-specific packets.  The open bridge is used only by a
separate export module after ordinary endpoints are stable.

\subsection{Structural kernel}

The accompanying \code{lean/SCOPE/Kernel.lean} records source-oriented wrappers
around checked rules:

\begin{lstlisting}
structure Contract (n : Nat) (Gamma : List (Fm PTm n))
    (goal : Fm PTm n) where
  valid : Fm PTm (n + 1)
  extract : valid :: Gamma.map Fm.shift
      |- [PRAAx] goal.shift

namespace Contract

theorem close (C : Contract n Gamma goal)
    (producer : Gamma |- [PRAAx] Fm.ex C.valid) :
    Gamma |- [PRAAx] goal :=
  .exE producer C.extract

theorem conditional (C : Contract n Gamma goal) :
    Gamma |- [PRAAx] (Fm.ex C.valid) --> goal := by
  apply Prf.impI
  apply Prf.exE Prf.hyp0
  exact C.extract.weaken (by simp)
end Contract
\end{lstlisting}

The file also wraps \code{Prf.trans}, \code{Prf.subst}, weakening, and boundary
discharge.  It contains no placeholder proof.  The first formalization
milestone is simply to elaborate it at the pinned revision and add it to the
axiom audit.

\subsection{Atomic schemas and bounded codes}

The schema-first interface is represented by:

\begin{lstlisting}
structure AtomSchema (k : Nat) where
  formula : Fm PTm k
  bit : PR k
  boolean : ([] : List (Fm PTm k)) |- [PRAAx]
    (app bit var == zero) \/ (app bit var == one)
  one_iff : ([] : List (Fm PTm k)) |- [PRAAx]
    Fm.iff (app bit var == one) formula

inductive BCode : Nat -> Type
  | atom  : AtomSchema k -> BCode k
  | conj  : BCode k -> BCode k -> BCode k
  | disj  : BCode k -> BCode k -> BCode k
  | impl  : BCode k -> BCode k -> BCode k
  | neg   : BCode k -> BCode k
  | exLt  : PTm k -> BCode (k + 1) -> BCode k
  | allLt : PTm k -> BCode (k + 1) -> BCode k
\end{lstlisting}

The actual archive uses repository notation in the definition of
\code{BCode.denote}.  An atom package stores only truth reflection; a generic
lemma derives falsity reflection from Booleanity.

\subsection{Arity-generic search}

The critical arithmetic implementation task is a parameterized version of
\code{Core/Search.lean}.  A suitable API is:

\begin{lstlisting}
namespace SearchN

def muF (body : PTm (k + 1)) : PR (k + 1)
def mu (body : PTm (k + 1)) (N : PTm n)
    (xs : Fin k -> PTm n) : PTm n

theorem search_spec (N i : PTm n) (xs : Fin k -> PTm n) :
  Gamma |- [PRAAx]
    le (mu body N xs) N /\
    ((lt (mu body N xs) N -->
       subst (Fin.cases (mu body N xs) xs) body == zero) /\
     (lt i (mu body N xs) -->
       neg (subst (Fin.cases i xs) body == zero)))

end SearchN
\end{lstlisting}

The exact vector expressions will use the repository's \code{Fin.cases} and
substitution conventions.  The proof should be obtained by parameterizing the
existing \code{spec\_open}; no new mathematical idea is required.  A useful
engineering strategy is to refactor the current two-argument search theorem as
the \(k=1\) corollary of \code{SearchN}.

\subsection{The reflection record}

The target is:

\begin{lstlisting}
structure Reflection {k : Nat} (A : BCode k) where
  char : PR k
  boolean : ([] : List (Fm PTm k)) |- [PRAAx]
    (app char var == zero) \/ (app char var == one)
  one_iff : ([] : List (Fm PTm k)) |- [PRAAx]
    Fm.iff (app char var == one) A.denote
  zero_iff : ([] : List (Fm PTm k)) |- [PRAAx]
    Fm.iff (app char var == zero) (neg A.denote)

def reflect (A : BCode k) : Reflection A
\end{lstlisting}

For atoms, return the stored symbol and proofs.  For connectives, build symbols
by \code{PR.comp} and proof terms by finite case analysis.  At bounded
quantifiers, construct a \code{SearchN.muF} symbol from the body characteristic
symbol and compose it with the bound's \code{PTm.toPR} symbol.

The substitution theorem should be almost definitional:

\begin{lstlisting}
@[simp] theorem char_subst (A : BCode k)
    (sigma : Fin k -> PTm n) :
  Syn.subst sigma (app (reflect A).char var) =
    app (reflect A).char sigma := by
  rfl
\end{lstlisting}

Depending on elaboration, the right proof may be \code{rfl} or the existing
\code{PTm.subst\_app} simplifier.  The main point is that no recursive theorem
about recompilation appears.

\subsection{Observer packages}

Exact observer erasure should be implemented first:

\begin{lstlisting}
structure ObserverContract (n : Nat) where
  valid : Fm PTm (n + 1)
  obs : PTm (n + 1)
  spec : Fm PTm (n + 1)
  extract : valid :: Gamma.map Fm.shift |- [PRAAx] spec
  coherent : valid.shift :: valid.reindexSecond ::
      Gamma.map (fun A => A.shift.shift) |- [PRAAx]
      obsFirst == obsSecond
\end{lstlisting}

The exact de Bruijn reindexing helpers should be defined once and tested on
small examples.  A public \code{out} formula existentially binds the
certificate and equates a public output variable with \code{obs}.  The module
then proves output existence and uniqueness.  Quotient observers replace
\code{==} by a reflected relation; metric observers add a rational distance and
triangle-inequality library.

\subsection{Indexed plan compiler}

Implement \code{Plan} only after the structural kernel and reflection module
compile.  Its first constructors are exact, specialize, supply, seal, and
boundary.  A simple structural function returns \code{Prf}.  The graph DSL is
strictly secondary: it should generate a \code{Plan}, not bypass it.

The effect checker is ordinary decidable host code.  It rejects:

\begin{itemize}
\item an unknown dependency or cycle not packaged as a trace node;
\item a missing producer with no named residual;
\item a relational output without a coherence derivation;
\item a named output without a direct producer or finite bound;
\item a sealed node whose certificate-dependent term is requested publicly.
\end{itemize}

Soundness of the checker is not trusted as an axiom.  Even if it accepts an
invalid plan, the invalid plan cannot produce a well-typed final \code{Prf}.

\subsection{Trace and demand modules}

A generic trace package stores fixed terms and local derivations:

\begin{lstlisting}
structure TraceSystem (k : Nat) where
  init : PTm k
  step : PTm (k + 2)
  inv : Fm PTm (k + 2)
  invQF : inv.QF
  base : ...
  next : ...
\end{lstlisting}

Use \code{recDef} for the state and \code{IPRA.ind} for the invariant theorem.
For capped demand closure, begin with fixed dimension at the metalevel.  Store
monotonicity as an object derivation, not as a semantic function property.  The
rank, inflation, stabilization, and leastness proofs are quantifier-free.

\subsection{Module layout}

\begin{verbatim}
IPRADerivations/Framework/SCOPE/
  Kernel.lean             -- structural wrappers
  Bits.lean               -- Boolean PR symbols and reflection
  SearchN.lean            -- arbitrary parameter arity
  BoundedCode.lean        -- schema, denotation, compiler
  Selection.lean          -- least accepted code
  Observer.lean           -- exact/quotient/metric output
  Trace.lean              -- deterministic traces
  Demand.lean             -- capped cyclic closure
  Viability.lean          -- whole finite paths
  Plan.lean               -- indexed proof plan and compiler
  Normalize.lean          -- host graph normalization
  Examples/
    LeastDivisor.lean
    CubicBisection.lean
    StableObserver.lean
    HeightList.lean
IPRABridge/SCOPE.lean      -- optional backend corollaries
\end{verbatim}

\subsection{Formalization order and tests}

The recommended order is:

\begin{enumerate}[label=\textbf{M\arabic*.}]
\item elaborate the structural kernel;
\item refactor arity-generic search;
\item prove bit combinators and atomic equality/order schemas;
\item implement bounded reflection and nested-quantifier tests;
\item add least selection;
\item add exact relational observers;
\item implement the indexed plan compiler;
\item add traces, demand closure, and viability;
\item migrate four pilot theorems;
\item expose bridge corollaries and run the repository axiom audit.
\end{enumerate}

Every final theorem module must contain no \code{sorry}, no theorem-specific
\code{axiom}, and a literal ordinary endpoint.  Tests should substitute
nontrivial compound terms under two alternating bounded binders; this is the
most direct regression test for the schema-first design.
